\documentclass[12pt]{article}

\usepackage{sbc-template}
\usepackage{graphicx,url}
\usepackage{float}
\usepackage{pgfplots}
\usepackage{pgfplotstable}
\usepackage{tikz}
\usepackage[utf8]{inputenc}
\usepackage[brazil]{babel}

\pgfplotsset{compat=1.18}

\sloppy

\title{Análise dos Diferentes Procedimentos de Substituição de Memória Cache 
para Determinados Algoritmos}
\author{Guilherme S. L. Lopes -- RA143630}
\address{Universidade Estadual de Maringá -- UEM}

\begin{document}

\maketitle

\begin{resumo}
Este trabalho estuda como diferentes políticas de substituição afetam o 
desempenho de um computador. Foram realizados testes com o simulador gem5 
utilizando seis políticas diferentes e três tipos de programas. Os resultados 
mostraram que a forma como o código é compilado torna a execução muito mais 
rápida, chegando a reduzir o tempo em até 85\%, mesmo que isso diminua a taxa 
de acertos na cache. Além disso, percebeu-se que aumentar a quantidade de dados 
processados melhorou levemente a eficiência da cache.
\end{resumo}

\section{Introdução}
A memória cache foi desenvolvida com o intuito de aumentar a eficiência de um 
computador, diminuindo a quantia necessária de buscas na memória principal 
transferindo um bloco de memória por vez, ao invés de apenas uma única palavra
\cite{stallings2017arquitetura}.

Tendo isso em mente, foram-se desenvolvidas ao longo do tempo diversas formas de 
manipular a cache, visto que quando ela se enche, é necessário fazer uma troca 
nos blocos que estão sendo alocados nela, essas técnicas são chamadas de 
políticas de substituição. Como existem inúmeros diferentes algoritmos, nem 
todas as políticas são úteis, as vezes inclusive sendo realizado mais buscas 
que o necessário, fazendo com que a cache tenha seu um efeito contrário à seu 
propósito, por isso é importante avaliar quando uma política é útil e quando 
ela não é.

Pensando na efetividade das políticas de substituição, foi-se feita uma pesquisa 
para averiguar, qual é mais eficiente em programas de multiplicação de matrizes 
consecutivas, de cálculo de implementação de propagação de ondas e de cálculo 
de difusão de calor num espaço 3D -- 3mm, fdtd-2d e heat-3d respectivamente 
\cite{polybench} -- a fim de averiguar qual melhor se aplica a esses três.

\section{Experimentos}
    \subsection{Metodologia}
    Para esse estudo foi utilizado o simulador gem5 \cite{Binkert:2011:gem5}, 
    onde foi configurado uma máquina virtual com um processador single-core 
    com velocidade de clock de 3GHz, uma RAM com frequencia de 1600MHz e 
    capacidade de memória de 8GB e, com níveis L1 e L2 da cache com tamanho de 
    32KiB. 

    Com essas configurações, foi feito testes usando o polybench nos algoritmos 
    supracitados, usando de entrada mini e small datasets, com os tipos de 
    compilação O0 e O3, sendo uma compilação padrão, sem otimizações, e uma 
    compilação com a maior quantia de otimização, respectivamente.

    No experimento, serão usados 6 políticas de substituição: FIFO, LFU, LRU, 
    MRU, RANDOM e WEIGHTED LRU. Dentro dessas, a única não vista em sala foi a 
    ultima citada, que consiste em uma variação da LRU padrão, onde é inserido 
    pesos nos blocos a serem retirados, o com o menor peso é evitado e quando 
    há dois com pesos iguais, evita-se o mais antigo \cite{article}.

    \subsection{Resultados}
    
    \begin{figure}[H]
    \centering
    \begin{tikzpicture}
    % --- Eixo base: O0 (barras largas, fundo) ---
    \begin{axis}[
        ybar,
        bar width=6pt,
        width=\columnwidth,
        height=6cm,
        ymin=75, ymax=100,
        ylabel={Taxa de hits (\%)},
        ylabel style={font=\small},
        xtick={1,2,3,4,5,6},
        xticklabels={3MM\\Mini, 3MM\\Small, HEAT-3D\\Mini,
                    HEAT-3D\\Small, FDTD-2D\\Mini, FDTD-2D\\Small},
        x tick label style={align=center, font=\scriptsize},
        y tick label style={font=\scriptsize},
        ymajorgrids=true,
        grid style={dashed, gray!30},
        enlarge x limits=0.12,
        legend style={
            at={(0.5,-0.32)}, anchor=north,
            legend columns=6, font=\scriptsize,
        },
        legend cell align={left},
    ]
    \addplot[fill=blue!30,   draw=blue!50] coordinates{(1,98.63)(2,98.85)(3,98.66)(4,99.00)(5,98.64)(6,98.65)};
    \addplot[fill=teal!30,   draw=teal!50] coordinates{(1,98.62)(2,98.84)(3,98.65)(4,98.99)(5,98.64)(6,98.64)};
    \addplot[fill=orange!30, draw=orange!55] coordinates{(1,98.63)(2,98.70)(3,98.66)(4,99.01)(5,98.64)(6,98.69)};
    \addplot[fill=violet!30, draw=violet!50] coordinates{(1,98.63)(2,98.85)(3,98.66)(4,99.00)(5,98.64)(6,98.65)};
    \addplot[fill=red!30,    draw=red!50] coordinates{(1,98.54)(2,98.53)(3,98.55)(4,98.79)(5,98.55)(6,98.54)};
    \addplot[fill=brown!30,  draw=brown!50] coordinates{(1,98.62)(2,98.85)(3,98.65)(4,98.99)(5,98.63)(6,98.65)};
    \legend{W-LRU O0, FIFO O0, LFU O0, LRU O0, MRU O0, RANDOM O0}
    \end{axis}

    % --- Eixo sobreposto: O3 (barras estreitas, frente) ---
    \begin{axis}[
        ybar,
        bar width=6pt,
        width=\columnwidth,
        height=6cm,
        ymin=75, ymax=100,
        ylabel={},
        xticklabels={,,},         % oculta xticks duplicados
        yticklabels={,,},         % oculta yticks duplicados
        axis x line=none,
        axis y line=none,
        enlarge x limits=0.12,
        legend style={
            at={(0.5,-0.45)}, anchor=north,
            legend columns=6, font=\scriptsize,
        },
        legend cell align={left},
    ]
    \addplot[fill=blue!80,   draw=blue!95]   coordinates{(1,87.70)(2,90.39)(3,87.93)(4,91.54)(5,87.96)(6,89.87)};
    \addplot[fill=teal!80,   draw=teal!95]   coordinates{(1,87.70)(2,90.37)(3,87.92)(4,91.54)(5,87.96)(6,89.87)};
    \addplot[fill=orange!80, draw=orange!95] coordinates{(1,87.70)(2,89.31)(3,87.92)(4,91.73)(5,87.94)(6,90.17)};
    \addplot[fill=violet!80, draw=violet!95] coordinates{(1,87.70)(2,90.39)(3,87.93)(4,91.54)(5,87.96)(6,89.87)};
    \addplot[fill=red!80,    draw=red!95]    coordinates{(1,87.68)(2,89.05)(3,87.82)(4,91.50)(5,87.90)(6,89.94)};
    \addplot[fill=brown!80,  draw=brown!95]  coordinates{(1,87.70)(2,90.46)(3,87.92)(4,91.55)(5,87.96)(6,89.97)};
    \legend{W-LRU O3, FIFO O3, LFU O3, LRU O3, MRU O3, RANDOM O3}
    \end{axis}

    \end{tikzpicture}
    \caption{Taxa de hits por política}
    \label{fig:hits_overlap}
    \end{figure}

    \begin{figure}[H]
    \centering
    \begin{tikzpicture}
    % --- O0: fundo ---
    \begin{axis}[
        ybar,
        bar width=6pt,
        width=\columnwidth,
        height=6cm,
        ymin=0, ymax=90,
        ylabel={Tempo de execução (s)},
        ylabel style={font=\small},
        xtick={1,2,3,4,5,6},
        xticklabels={3MM\\Mini, 3MM\\Small, HEAT-3D\\Mini,
                    HEAT-3D\\Small, FDTD-2D\\Mini, FDTD-2D\\Small},
        x tick label style={align=center, font=\scriptsize},
        y tick label style={font=\scriptsize},
        ymajorgrids=true,
        grid style={dashed, gray!30},
        enlarge x limits=0.12,
        legend style={
            at={(0.5,-0.32)}, anchor=north,
            legend columns=6, font=\scriptsize,
        },
        legend cell align={left},
    ]
    \addplot[fill=blue!30,   draw=blue!50] coordinates{(1,30.07)(2,38.23)(3,30.07)(4,76.58)(5,28.92)(6,44.02)};
    \addplot[fill=teal!30,   draw=teal!50] coordinates{(1,28.67)(2,38.37)(3,30.92)(4,75.36)(5,28.92)(6,43.57)};
    \addplot[fill=orange!30, draw=orange!55] coordinates{(1,28.85)(2,38.64)(3,30.32)(4,75.35)(5,29.31)(6,43.59)};
    \addplot[fill=violet!30, draw=violet!50] coordinates{(1,28.76)(2,38.48)(3,30.11)(4,75.57)(5,29.17)(6,44.33)};
    \addplot[fill=red!30,    draw=red!50] coordinates{(1,29.01)(2,38.79)(3,30.47)(4,75.85)(5,29.22)(6,43.99)};
    \addplot[fill=brown!30,  draw=brown!50] coordinates{(1,28.78)(2,38.52)(3,30.01)(4,75.05)(5,28.90)(6,43.48)};
    \legend{W-LRU O0, FIFO O0, LFU O0, LRU O0, MRU O0, RANDOM O0}
    \end{axis}

    % --- O3: frente ---
    \begin{axis}[
        ybar,
        bar width=6pt,
        width=\columnwidth,
        height=6cm,
        ymin=0, ymax=90,
        ylabel={},
        xticklabels={,,},
        yticklabels={,,},
        axis x line=none,
        axis y line=none,
        enlarge x limits=0.12,
        legend style={
            at={(0.5,-0.45)}, anchor=north,
            legend columns=6, font=\scriptsize,
        },
        legend cell align={left},
    ]
    \addplot[fill=blue!80,   draw=blue!95]   coordinates{(1,6.98)(2,9.06)(3,7.25)(4,10.27)(5,7.03)(6,9.15)};
    \addplot[fill=teal!80,   draw=teal!95]   coordinates{(1,6.95)(2,8.99)(3,7.12)(4,10.15)(5,6.95)(6,9.03)};
    \addplot[fill=orange!80, draw=orange!95] coordinates{(1,7.00)(2,9.19)(3,7.27)(4,10.21)(5,7.01)(6,9.08)};
    \addplot[fill=violet!80, draw=violet!95] coordinates{(1,7.01)(2,9.05)(3,7.19)(4,10.24)(5,7.03)(6,9.18)};
    \addplot[fill=red!80,    draw=red!95]    coordinates{(1,7.04)(2,9.01)(3,7.10)(4,10.27)(5,7.02)(6,9.14)};
    \addplot[fill=brown!80,  draw=brown!95]  coordinates{(1,6.93)(2,8.94)(3,7.10)(4,10.14)(5,6.97)(6,8.99)};
    \legend{W-LRU O3, FIFO O3, LFU O3, LRU O3, MRU O3, RANDOM O3}
    \end{axis}

    \end{tikzpicture}
    \caption{Tempo de execução por política}
    \label{fig:exec_overlap}
    \end{figure}

    Como é possível observar na figura \ref{fig:hits_overlap}, usando o método 
    de compilação com mais acertos de cache, foi justamente aquele com o menor 
    nível de otimização, mantendo uma média de hits próxima dos 98\%, porém 
    também aumentou exponencialmente o tempo necessário para a execução, visível 
    na figura \ref{fig:exec_overlap}. 
    
    É interessante de se notar especialmente que com a compilação O3 que, por mais 
    que teve uma taxa de hits menor, então era de se esperar um tempo de execução 
    maior, devido a otimização na leitura dos algoritmos pelo compilador, o 
    tempo de execução é no minimo menor que a metade, e no caso do heat-3d pra 
    um dataset small, é aproximadamente 15\% do tempo de quando usamos a 
    compilação O0. 

    Outra possível observação é que para o dataset de tamanho small, a taxa de 
    hits -- mesmo que não por muito -- foi maior, aumentando em cerca de 2 à 3\%, 
    com 3 segundos a mais na execução em média. Como os algoritmos escolhidos 
    são algoritmos mais complexos, com multiplicações continuas de matrizes e 
    cálculos espaciais, a expectativa era que os o tempo de execução para o 
    usando mais entradas seria obviamente maior, porém também é interessante 
    notar que os hits também foram maiores, tendo aproximadamente um aumento de 
    12 milhões de hits, com 150 mil misses a mais, excluindo a MRU, que teve em 
    todos um aumento expressivo de misses, alcançando 300 mil a mais na média.

\section{Conclusão}
Analisando os dados, podemos ver que a taxa de hits na compilação O0 é melhor em 
todos os algortimos, porém o elevado tempo de execução torna inviável dependendo 
da quantia de vezes que é necessária rodá-lo. Um possível cenário onde possa ser 
mais interessante o uso desse modo, é com algoritmos menores e menos complexos, 
visto que o tempo necessário reduziria e os acertos da cache se manteriam elevados.

Já para a compilação O3, mesmo mantendo uma quantia menores de acertos, o tempo 
de execução se reduziu tanto a ponto de talvez ser mais viável utilizar esse 
método para algoritmos complexos, sacrificando a taxa de hits para ganhar mais 
execuções. Possivelmente, com um dataset maior ainda os acertos podem melhorar, 
uma vez que teve um aumento no dataset mini para o dataset small, é possível 
considerar uma outra pesquisa para verificar se com datasets maiores ainda, 
essa ideia ainda é válida.   
    

\bibliographystyle{sbc}
\bibliography{relatorio}

\end{document}